% CC BY 4.0 License
% Copyright (c) 2026 Nanum Learn

\documentclass[a4paper,11pt]{article}

\usepackage[T1]{fontenc}
\usepackage[utf8]{inputenc}
\usepackage[english]{babel}

\usepackage{amsmath}
\usepackage{amssymb}
\usepackage{amsthm}
\usepackage{mathtools}

\usepackage[margin=1in]{geometry}
\usepackage{lastpage}
\usepackage{fancyhdr}
\fancyhead[L]{NL01CHEM01 Session 2 Problem Set Solutions}
\fancyhead[R]{Nanum Learn}
\fancyfoot[C]{}
\pagestyle{fancy}
\usepackage[skip=10pt]{parskip}
\setlength{\headheight}{13.59999pt}
\setlength{\headsep}{8pt}

\usepackage{titlesec}
\titleformat{\section}
    {\normalfont\LARGE\bfseries\color{BrickRed}}
    {\S\! \thesection}
    {0.5em}
    {}
\titleformat{\subsection}
    {\normalfont\Large\bfseries}
    {\thesubsection}
    {0.5em}
    {}

\usepackage{enumitem}
\usepackage[dvipsnames]{xcolor}


\newcommand{\easy}{(\ensuremath{\bullet \circ \circ})}
\newcommand{\medium}{(\ensuremath{\bullet \bullet \circ})}
\newcommand{\hard}{(\ensuremath{\bullet \bullet \bullet})}

\newtheoremstyle{listprobstyle}
    {\parsep}{\parsep}
    {\normalfont}
    {}
    {\bfseries}
    {}
    { }
    {{\thmname{#1}~\thmnumber{#2}.}~
        \thmnote{\textcolor{BrickRed}{\bfseries #3}}\\}
\theoremstyle{listprobstyle}
\newtheorem{problem}{Problem}
\newenvironment{solution}
    {\par\begingroup\color{BrickRed}}
    {\par\endgroup}



\title{\textbf{NL01CHEM01 Session 2 Problem Set Solutions}}
\date{10 November 2026}
\author{William Kang}

\begin{document}

\maketitle
\begin{center}
    \small\textit{Edited and Formatted by Enoch Yu}
\end{center}

\section{Solutions}
\begin{problem}[Outer Electrons from Column \easy]
How many outer electrons do elements in each column have?
\vspace{-0.6\baselineskip}
\begin{enumerate}[label=\alph*),itemsep=0pt,parsep=0pt,topsep=0pt]
\item Column 1 (lithium family):
    \underline{\textcolor{BrickRed}{\phantom{X}1\phantom{X}}}
\vspace{0.5\baselineskip}
\item Column 14 (carbon family):
    \underline{\textcolor{BrickRed}{\phantom{X}4\phantom{X}}}
\vspace{0.5\baselineskip}
\item Column 17 (chlorine family):
    \underline{\textcolor{BrickRed}{\phantom{X}7\phantom{X}}}
\end{enumerate}
\end{problem}

\begin{problem}[Column from Outer Count \easy]
Identify the periodic table column from the number of outer electrons.
\vspace{-0.6\baselineskip}
\begin{enumerate}[label=\alph*),itemsep=0pt,parsep=0pt,topsep=0pt]
\item 2 outer electrons:
    Column
    \underline{\textcolor{BrickRed}{\phantom{X}2\phantom{X}}}
\vspace{0.5\baselineskip}
\item 6 outer electrons:
    Column
    \underline{\textcolor{BrickRed}{\phantom{X}16\phantom{X}}}
\vspace{0.5\baselineskip}
\item 7 outer electrons:
    Column
    \underline{\textcolor{BrickRed}{\phantom{X}17\phantom{X}}}
\end{enumerate}
\end{problem}

\begin{problem}[Write the Ion \easy]
For each atom, write the ion it usually forms.
\vspace{-0.6\baselineskip}
\begin{enumerate}[label=\alph*),itemsep=0pt,parsep=0pt,topsep=0pt]
\item Sodium: \underline{
        \textcolor{BrickRed}{\phantom{X}$\mathrm{Na^+}$\phantom{X}}
    }
    \begin{solution}
    Sodium loses 1 electron.
    \end{solution}
\vspace{0.5\baselineskip}
\item Magnesium: \underline{
        \textcolor{BrickRed}{\phantom{X}$\mathrm{Mg^{2+}}\phantom{X}$}
    }
    \begin{solution}
    Magnesium loses 2 electrons.
    \end{solution}
\vspace{0.5\baselineskip}
\item Chlorine: \underline{
        \textcolor{BrickRed}{\phantom{X}$\mathrm{Cl^-}$\phantom{X}}
    }
    \begin{solution}
    Chlorine gains 1 electron.
    \end{solution}
\vspace{0.5\baselineskip}
\item Oxygen: \underline{
        \textcolor{BrickRed}{\phantom{X}$\mathrm{O^{2-}}\phantom{X}$}
    }
    \begin{solution}
    Oxygen gains 2 electrons.
    \end{solution}
\end{enumerate}
\end{problem}


\begin{problem}[Ion or Share? \easy]
Label each compound I (ionic) or C (covalent).
\vspace{-0.6\baselineskip}
\begin{enumerate}[label=\alph*),itemsep=0pt,parsep=0pt,topsep=0pt]
\item Salt ($\mathrm{NaCl}$): (\textcolor{BrickRed}{I})
\begin{solution}
Sodium transfers an electron to chlorine and the resulting oppositely
charged ions attract.
\end{solution}
\vspace{0.5\baselineskip}

\item Water ($\mathrm{H_2O}$): (\textcolor{BrickRed}{C})
\begin{solution}
Hydrogen and oxygen share electrons, forming covalent bonds.
\end{solution}
\vspace{0.5\baselineskip}

\item $\mathrm{MgCl_2}$: (\textcolor{BrickRed}{I})
\begin{solution}
Magnesium forms $\mathrm{Mg^{2+}}$ and chlorine forms $\mathrm{Cl^-}$.
\end{solution}
\vspace{0.5\baselineskip}

\item $\mathrm{CO_2}$: (\textcolor{BrickRed}{C})
\begin{solution}
Carbon and oxygen share electrons and $\mathrm{CO_2}$ is covalent.
\end{solution}
\end{enumerate}
\end{problem}


\begin{problem}[Write the Formula \easy]
Write the neutral chemical formula formed by each pair of ions.
\vspace{-0.6\baselineskip}
\begin{enumerate}[label=\alph*),itemsep=0pt,parsep=0pt,topsep=0pt]
\item Sodium ($\mathrm{Na^+}$) and chlorine ($\mathrm{Cl^-}$):
    \underline{
        \textcolor{BrickRed}{\phantom{X}$\mathrm{NaCl}$\phantom{X}}
    }
    \begin{solution}
    Sodium has charge $+1$ and chlorine has charge $-1$.
    \end{solution}
\vspace{0.5\baselineskip}

\item Magnesium ($\mathrm{Mg^{2+}}$) and oxygen ($\mathrm{O^{2-}}$):
    \underline{
        \textcolor{BrickRed}{\phantom{X}$\mathrm{MgO}$\phantom{X}}
    }
    \begin{solution}
    Magnesium has charge $+2$ and oxygen has charge $-2$.
    \end{solution}
\end{enumerate}
\end{problem}


\begin{problem}[Count Atoms in a Formula \easy]
Consider the chemical formula $\mathrm{Al_2O_3}$.
\vspace{-0.6\baselineskip}
\begin{enumerate}[label=\alph*),itemsep=0pt,parsep=0pt,topsep=0pt]
\item How many aluminum atoms?
    \underline{\textcolor{BrickRed}{\phantom{X}2\phantom{X}}}
\vspace{0.5\baselineskip}
\item How many oxygen atoms?
    \underline{\textcolor{BrickRed}{\phantom{X}3\phantom{X}}}
\vspace{0.5\baselineskip}
\item What is the Al:O ratio?
    \underline{\textcolor{BrickRed}{\phantom{X}$2:3$\phantom{X}}}
\end{enumerate}
\end{problem}


\begin{problem}[Lose or Gain? \easy]
For each ion, state whether the atom lost or gained electrons.
\vspace{-0.6\baselineskip}
\begin{enumerate}[label=\alph*),itemsep=0pt,parsep=0pt,topsep=0pt]
\item $\mathrm{Na^+}$:
    \underline{\textcolor{BrickRed}{Lost}}
    \begin{solution}
    Sodium loses 1 electron to form $\mathrm{Na^+}$.
    \end{solution}
\vspace{0.5\baselineskip}
\item $\mathrm{Cl^-}$:
    \underline{\textcolor{BrickRed}{Gained}}
    \begin{solution}
    Chlorine gains 1 electron to form $\mathrm{Cl^-}$.
    \end{solution}
\vspace{0.5\baselineskip}
\item $\mathrm{Mg^{2+}}$:
    \underline{\textcolor{BrickRed}{Lost 2}}
    \begin{solution}
    Magnesium loses 2 electrons to form $\mathrm{Mg^{2+}}$.
    \end{solution}
\vspace{0.5\baselineskip}
\item $\mathrm{O^{2-}}$:
    \underline{\textcolor{BrickRed}{Gained 2}}
    \begin{solution}
    Oxygen gains 2 electrons to form $\mathrm{O^{2-}}$.
\end{solution}
\end{enumerate}
\end{problem}


\begin{problem}[True or False? \easy]
Write T (true) or F (false) next to each statement.
\vspace{-0.6\baselineskip}
\begin{enumerate}[label=\alph*),itemsep=0pt,parsep=0pt,topsep=0pt]
\item $\mathrm{MgCl_2}$ is neutral overall.
    (\textcolor{BrickRed}{T})
    \begin{solution}
    One $\mathrm{Mg^{2+}}$ ion and two $\mathrm{Cl^-}$
    ions give total charge $2 - 2=0$.
    \end{solution}
\vspace{0.5\baselineskip}
\item $\mathrm{Na^+}$ has more electrons than protons.
    (\textcolor{BrickRed}{F})
    \begin{solution}
    $\mathrm{Na^+}$ has lost 1 electron.
    \end{solution}
\vspace{0.5\baselineskip}
\item Subscripts in chemical formulas indicate charge.
    (\textcolor{BrickRed}{F})
    \begin{solution}
    Subscripts count atoms or ions in a chemical formula and
    superscripts indicate ionic charge.
    \end{solution}
\vspace{0.5\baselineskip}
\item Elements in the same column have the same number of electron
    shells.
    (\textcolor{BrickRed}{F})
    \begin{solution}
    Elements in the same row (period) have the same number of occupied
    electron shells and columns generally share the same number of
    valence electrons.
    \end{solution}
\end{enumerate}
\end{problem}


\begin{problem}[A Trickier Pair \medium]
Calcium forms a $2+$ ion ($\mathrm{Ca^{2+}}$), and hydroxide
is a $1-$ ion ($\mathrm{OH^-}$).
\vspace{-0.6\baselineskip}
\begin{enumerate}[label=\alph*),itemsep=0pt,parsep=0pt,topsep=0pt]
\item Write the neutral chemical formula:
    \underline{\textcolor{BrickRed}{$\mathrm{Ca(OH)_2}$}}
    \begin{solution}
    One calcium ion has charge $+2$ and two hydroxide ions each have
    charge $-1$.
    \end{solution}
\vspace{0.5\baselineskip}
\item How many calcium ions are needed?
    \underline{\textcolor{BrickRed}{1}}
    \begin{solution}
    One $\mathrm{Ca^{2+}}$ ion provides charge $+2$.
    \end{solution}
\vspace{0.5\baselineskip}
\item How many hydroxide ions are needed?
    \underline{\textcolor{BrickRed}{2}}
    \begin{solution}
    Two $\mathrm{OH^-}$ ions provide charge $-2$.
    \end{solution}
\end{enumerate}
\end{problem}


\begin{problem}[Water in Lewis Form \medium]
Consider the Lewis structure of water ($\mathrm{H_2O}$).
\vspace{-0.6\baselineskip}
\begin{enumerate}[label=\alph*),itemsep=0pt,parsep=0pt,topsep=0pt]
\item How many shared electron pairs are there?
    \underline{\textcolor{BrickRed}{\phantom{X}2\phantom{X}}}
    \begin{solution}
    Oxygen forms two covalent bonds and each bond contains one shared
    electron pair.
    \end{solution}
\vspace{0.5\baselineskip}
\item How many lone pairs are on the oxygen atom?
    \underline{\textcolor{BrickRed}{\phantom{X}2\phantom{X}}}
    \begin{solution}
    Oxygen has two lone pairs that are not involved in bonding.
    \end{solution}
\vspace{0.5\baselineskip}
\item How many electrons surround the oxygen atom?
    \underline{\textcolor{BrickRed}{\phantom{X}8\phantom{X}}}
    \begin{solution}
    Oxygen has 2 shared pairs and 2 lone pairs where each pair contains
    2 electrons. Therefore, there are $(2 + 2) \cdot 2 = 8$ electrons
    around oxygen.
    \end{solution}
\end{enumerate}
\end{problem}
\end{document}


